% !TeX root = ../main.tex
\section{AI Đồng Hành}
\label{sec:ai}

AI đồng hành cho phép người chơi luyện tập khi không có đối thủ.

\subsection{Vì Sao Không Dùng MCTS}

Phần mềm không dùng Monte Carlo Tree Search (MCTS)~\citep{browne2012survey}.
MCTS cổ điển đánh giá một nước đi bằng cách chơi thử ngẫu nhiên đến hết ván
rất nhiều lần, mà ván Cờ Trận thì rất dài: luật không có giới hạn số nước,
và một ván chơi ngẫu nhiên kéo dài trung vị khoảng 300 nước ở chế độ 2
người, 1.300--1.600 nước ở chế độ 3 người và 3.300--3.500 nước ở chế độ 4
người (Mục~\ref{sec:eval_branching}). Trên cùng máy đo ở
Mục~\ref{sec:eval}, một lần chơi thử như vậy trong Python thuần mất trung vị
14\,ms, 106\,ms và 299\,ms (30 ván mỗi chế độ, tệp
\texttt{experiments/results/playout\_cost.csv}). Chỉ 1.000 lần chơi thử cho
một nước đi đã tốn khoảng 14 giây, gần 2 phút và gần 5 phút tương ứng, trong khi bot Minimax
hiện tại chỉ cần 0,17--3 giây mỗi nước (Bảng~\ref{tab:timing}). Ngược lại,
luật kẹp--lật cho sẵn một tín hiệu đánh giá rẻ và có ý nghĩa là quân số và
số nước đi còn lại, theo tinh thần các hàm đánh giá cho
Othello~\citep{buro1999how}. Vì vậy tìm kiếm Minimax có giới hạn độ sâu là
lựa chọn tự nhiên. Chúng tôi không so sánh thực nghiệm với MCTS và không
khẳng định thuật toán nào mạnh hơn.

\subsection{Tìm Kiếm Đối Kháng}
\label{ssec:search}

Bot dùng \emph{hai} thuật toán tách biệt tuỳ số người chơi---cùng họ
Minimax với cắt tỉa Alpha-Beta nhưng khác nhau ở hàm đánh giá và cách xử
lý kết cuộc. Ký hiệu $N_c$ là số quân và $M_c$ là số nước đi hợp lệ
(mobility) của phe $c$.

\paragraph{Chế độ 2 người---Minimax Alpha-Beta.} Hai phe đối kháng hoàn
toàn nên bot dùng Minimax cổ điển với cắt tỉa
Alpha-Beta~\citep{knuth1975alphabeta}: phe mình cực đại, đối thủ cực tiểu, luân
phiên qua cờ \textsc{max}. Hàm đánh giá kết hợp hai mục tiêu, ăn quân và hạn chế
số nước đi của đối thủ, bằng cách chuẩn hoá cả hai
số hạng về $[-1,1]$ rồi gán trọng số $0.7$ và $0.3$
(xem lý do ở Mục~\ref{ssec:pitfalls}):
\begin{equation}
\mathrm{eval}_2(b,c)=0.7\,\frac{N_c-N_{\bar c}}{N_c+N_{\bar c}}
+0.3\,\frac{M_c-M_{\bar c}}{M_c+M_{\bar c}}\ \in[-1,1].
\label{eq:eval}
\end{equation}
Khi một phe không còn sống, điểm kết cuộc kèm thưởng thắng nhanh
\begin{equation}
\mathrm{score}^{2}_{\text{term}}=\pm(10000+d)
\label{eq:term2}
\end{equation}
($+$ nếu đối thủ thua, $-$ nếu mình thua; $d$ là độ sâu còn lại). Giá trị
kết cuộc lớn hơn hẳn mọi giá trị đánh giá trong $[-1,1]$, nên việc thu nhỏ thang đo không
làm bot bỏ lỡ nước thắng/thua thực sự.

\paragraph{Chế độ 3--4 người---Paranoid Minimax.} Với $\ge 3$ phe, điểm
số không còn đối nghịch hoàn toàn nên cắt tỉa khó
hơn~\citep{korf1991multiplayer,sturtevant2000pruning}. Bot áp dụng
\textbf{Paranoid Minimax}~\citep{sturtevant2000pruning,sturtevant2003comparison}:
phe đang xét cực đại, \emph{mọi phe còn
lại bị coi như một liên minh} cùng cực tiểu, quy bài toán $N$-người về
max--min hai chiều để dùng Alpha-Beta đầy đủ. Giữ nguyên trọng số $0.7/0.3$
và số hạng quân số chuẩn hoá về $[-1,1]$ như chế độ 2 người, nhưng vì các
đối thủ bị gộp thành một liên minh, số hạng mobility được làm \emph{bất đối
xứng}: chỉ đo mobility $M_c$ của riêng phe đang xét (không trừ mobility đối
thủ) và chuẩn hoá theo tổng quân $T$ để giữ trong $[0,1)$:
\begin{equation}
\mathrm{eval}_{\text{par}}(b,c)=0.7\,\frac{N_c-N_{\bar c}}{T}
+0.3\,\frac{M_c}{M_c+T},\qquad T=N_c+N_{\bar c},
\label{eq:evalpar}
\end{equation}
với $N_{\bar c}$ là tổng quân của \emph{tất cả} đối thủ. Khác với
\eqref{eq:eval}, số hạng mobility ở đây không bao giờ âm và không so trực
tiếp với đối thủ; nó chỉ thưởng phe còn nhiều nước đi (tránh bị phong
toả)---một lựa chọn thực dụng vì mobility của một liên minh gộp không phản
ánh rõ mối đe doạ từ từng đối thủ riêng lẻ. Khi chỉ còn một
phe sống, điểm kết cuộc kèm thưởng thắng nhanh
\begin{equation}
\mathrm{score}^{\text{par}}_{\text{term}}=\pm(10000+d),
\label{eq:termpar}
\end{equation}
trong đó $d$ là độ sâu còn lại---khuyến khích bot kết liễu sớm thay vì
kéo dài ván. Trong cây tìm kiếm, phe đến lượt mà hết nước đi bị loại và
lượt chuyển sang phe kế tiếp, đúng như luật chơi. Một lựa chọn thay thế mạnh là Best Reply
Search~\citep{schadd2011brs}: nó vượt Paranoid trong Chinese Checkers và
Focus, nhưng chỉ ngang Paranoid trong Rolit~\citep{schadd2011brs,nijssen2013mcts}.
Chúng tôi chọn Paranoid vì đơn giản và đủ tốt trong ngân sách
tính toán thời gian thực, để dành so sánh với Best Reply Search cho hướng
phát triển.

\begin{table}[H]
  \centering
  \caption{Tham số AI theo chế độ chơi, đúng như trong phần mềm. Thời gian
  suy nghĩ đo thực tế ở Bảng~\ref{tab:timing}.}
  \label{tab:ai_params}
  \begin{tabular}{lcc}
    \toprule
    \textbf{Chế độ} & \textbf{Thuật toán} & \textbf{Độ sâu} \\
    \midrule
    2 người & Minimax Alpha-Beta & 6 \\
    3 người & Paranoid Minimax   & 4 \\
    4 người & Paranoid Minimax   & 3 \\
    \bottomrule
  \end{tabular}
\end{table}

Ở chế độ 3--4 người, AI chạy trong một luồng nền
(\texttt{asyncio.run\_in\_executor}) nên giao diện không bị đứng khi bot
suy nghĩ. Ở chế độ 2 người, bot hiện chạy trên luồng chính, nên giao diện
tạm dừng trong lúc bot tính nước.

\subsection{Hiệu quả cắt tỉa và sắp thứ tự nước đi}
\label{ssec:ordering}

Hiệu quả của cắt tỉa Alpha-Beta phụ thuộc mạnh vào thứ tự duyệt nước
đi~\citep{knuth1975alphabeta}.

Cờ Trận cho một heuristic sắp xếp rẻ và tự nhiên: vì mỗi hành
động \emph{thấy trước} số quân ăn được $|a|=|F_\delta|$
(Mục~\ref{ssec:capture}), bot duyệt các nước theo $|a|$ \emph{giảm dần}; riêng engine 2 người
còn duyệt trước tiên nước tốt nhất đã lưu trong bảng chuyển vị (xem dưới).
Nước ăn nhiều quân thường là ứng viên tốt, nên cắt tỉa xảy ra sớm hơn. Để
tăng tốc, cả hai engine dùng \emph{cập nhật/hoàn-tác nước} (make/unmake) thay cho việc sao
chép trạng thái ở mỗi nút. Riêng engine 2 người còn dùng một \textbf{bảng
băm chuyển vị} (transposition table) khoá theo băm Zobrist---lưu cận
trên/dưới/đúng theo Alpha-Beta và ưu tiên duyệt nước tốt nhất đã ghi---để
tránh tính lại các dãy nước dẫn về cùng một trạng thái; engine Paranoid
chưa dùng bảng chuyển vị vì khoá trạng thái (kèm tập phe còn sống và con
trỏ lượt) phức tạp hơn, để dành cho bản sau. Tìm kiếm tăng dần độ sâu
(iterative deepening) chưa được cài đặt và để dành cho hướng phát triển.

\subsection{Ba nhược điểm thiết kế và cách khắc phục}
\label{ssec:pitfalls}

Khi thử nghiệm bot, chúng tôi phát hiện ba nhược điểm \emph{không} phải lỗi
cài đặt mà bắt nguồn từ việc áp một thuật toán tổng quát lên một luật chơi
đặc thù: quân được bảo toàn (ăn $=$ đổi phe, Mục~\ref{ssec:capture}) và
điều kiện thắng thực chất là \emph{phong toả} đối thủ chứ không phải tiêu
diệt. Mỗi nhược điểm phản ánh một giả định ngầm sai trong sự khớp đó. Cả
ba đều có cách khắc phục cụ thể áp dụng cho cả hai engine (2 người và
Paranoid).

\paragraph{Nhược điểm \#1 --- Tham ăn cực đại (greedy capture).}
\emph{Nguyên nhân.} Hàm sinh nước đi mặc định, khi một nước tạo nhiều
hướng kẹp, chỉ giữ hướng lật \emph{nhiều quân nhất} ($\arg\max_\delta
|F_\delta|$). Bot không bao giờ xét nước ăn ít hơn, kể cả khi nước ăn to
mở ra phản đòn kẹp lại. Vì quân bảo toàn, một nước ăn 5 quân có thể bị
đối thủ lật lại $\ge 6$ quân ngay lượt sau---đúng kiểu ``nhử mồi'' mà luật
kẹp cho phép.
\emph{Khắc phục.} Thay tập hành động $\mathcal A(c,s)$ bằng tập \emph{mở
rộng} $\mathcal A^{+}$: mỗi ô đáp có $k\ge 2$ hướng kẹp sinh ra $k$ hành
động riêng, mỗi hành động ứng với một hướng. Bot do đó cân nhắc cả nước
ăn ít; chính cây tìm kiếm (qua eval ở các lớp sau) quyết định ăn bao nhiêu
là có lợi, thay vì một quy tắc tham lam cố định.

\paragraph{Nhược điểm \#2 --- Sai trọng tâm hàm đánh giá.}
\emph{Nguyên nhân.} Hàm đánh giá ban đầu, $10(N_c-N_{\bar c})+(M_c-M_{\bar
c})$, đặt trọng số quân số gấp mười lần mobility. Nhưng vì quân được bảo toàn,
ăn hết quân đối thủ rất khó; ván cờ thực tế kết thúc khi một phe bị phong
toả. Trong 100 ván giữa hai bot độ sâu~2 (Mục~\ref{sec:eval_tournament}),
không ván nào kết thúc vì bên thua hết quân: cả 100 ván kết thúc khi bên
thua bị phong toả trong lúc vẫn còn trung vị 8 quân. Bot vì thế tối ưu sai
mục tiêu---gom quân thay vì phong toả.
Ngoài ra hai số hạng khác đơn vị (đếm quân so với đếm nước đi) nên tỉ số
trọng số $10\!:\!1$ không mang ý nghĩa rõ ràng.
\emph{Khắc phục.} Chuẩn hoá cả hai số hạng về $[-1,1]$ rồi gán trọng số
phản ánh trực giác người chơi: $0.7$ cho ăn quân, $0.3$ cho chặn nước đi
đối thủ (Phương trình~\eqref{eq:eval} và~\eqref{eq:evalpar}). Trọng số của số
nước đi nhỏ hơn, nhưng vì hai số hạng \emph{cùng thang đo} nên số hạng này
vẫn ảnh hưởng thật sự tới lựa chọn của bot.

\paragraph{Nhược điểm \#3 --- Tất định và lặp nước.}
\emph{Nguyên nhân (hai tầng).} (i)~\emph{Tất định}: phép $\arg\max$ luôn
trả nước tốt nhất đầu tiên, nên với cùng một thế cờ bot luôn đi y hệt---tìm
được một đường thắng là thắng lại vô hạn. (ii)~\emph{Dao động}: vì eval là
số thực liên tục nên hầu như không bao giờ có hai nước bằng điểm tuyệt
đối, khiến chiến lược ``ngẫu nhiên khi hoà điểm'' vô hiệu. Tệ hơn, do quân
bảo toàn, bot thường đẩy một quân đi rồi kéo về đúng thế cũ
($A\!\to\!B\!\to\!A$), tạo vòng lặp. Đo trên 40 ván tự đánh ở độ sâu~2,
mỗi ván quan sát tối đa 200 nước đầu vì khi chưa sửa, bot có thể đi vòng
mãi không kết thúc (tổng 5040 lượt đi), \textbf{37{,}3\%} số thế cờ là một thế \emph{đã từng}
xuất hiện trong ván---cho thấy vòng dao động này phổ biến đến mức nào khi
không can thiệp. Điểm mấu chốt:
nhược điểm này \emph{không} sửa được bằng can thiệp vào tìm kiếm---Alpha-Beta,
sắp thứ tự nước, hay bảng chuyển vị chỉ đổi \emph{tốc độ} chứ không đổi
\emph{nước được chọn}---nên phải sửa ở khâu chọn nước.
\emph{Khắc phục (hai cơ chế).}
\begin{itemize}
  \item \textbf{Phạt lặp thế (anti-repetition).} Ván cờ lưu một lịch sử
        $\mathcal H$ gồm băm Zobrist của \emph{mọi} thế đã thực sự xảy ra.
        Ở gốc cây, một nước dẫn tới thế $b'$ đã thuộc $\mathcal H$ bị trừ
        điểm tỉ lệ với số lần lặp:
        $\mathrm{eval}'(a)=\mathrm{eval}(a)-\lambda\cdot\#\{b'\in\mathcal
        H\}$ với $\lambda=0.5$. Vì $\lambda$ lớn hơn nhiều khoảng cách eval
        giữa các nước hợp lý, bot từ bỏ vòng dao động, nhưng vẫn chấp nhận
        lặp nếu đó là nước duy nhất không dẫn tới thua.
  \item \textbf{Ngẫu nhiên trong biên $\varepsilon$.} Thay cho ``hoà điểm
        tuyệt đối'', mọi nước nằm trong biên $\varepsilon=0.02$ quanh nước
        tốt nhất được coi là tương đương và chọn ngẫu nhiên đều. Điều này
        tạo đa dạng giữa các ván mà không hi sinh chất lượng---các nước
        kém hơn quá $\varepsilon$ vẫn bị loại.
\end{itemize}
Sau khắc phục, lặp lại đúng thí nghiệm 40 ván nói trên cho tỉ lệ lặp thế
\textbf{0{,}0\%} (0 thế lặp trên 4069 lượt đi)---vòng dao động bị loại bỏ
hoàn toàn. Lưu ý lịch sử $\mathcal H$ chỉ chứa các thế \emph{đã} xảy ra
nên không làm tăng chi phí tìm kiếm theo cấp số nhân; nó chỉ là một tra cứu
băm $O(1)$ tại gốc.
